\chapter{Construcción analítica de los funcionales de la torre triádica}
\label{chap:funcionales-espectrales-triadicos}

La síntesis anterior ha situado las constantes espectrales dentro de la
torre triádica. Este bloque contiene su construcción analítica primaria:
sucesión de operadores autoadjuntos sobre ciclos de orden \(3^m\), límite de
trazas de calor y transformada de Mellin. Las constantes clásicas aparecen
como evaluaciones, partes finitas o derivadas del funcional meromorfo
resultante, no como prefijos reconocidos en un catálogo.

La normalización del operador $D_m$ contiene explícitamente $\pi$ en la escala
que relaciona el laplaciano discreto con el núcleo de calor gaussiano. El
apartado constituye, por tanto, una extensión analítica de la rama
arquimediana ya construida. Las constantes espectrales posteriores se obtienen
simultáneamente del funcional resultante, con esa normalización como hipótesis
de partida.

\section{Límites inductivos y escala de refinamiento}

La prolongación espectral exige distinguir dos niveles analíticos. El
primero corresponde a familias acotadas y compatibles; el segundo, a
operadores no acotados, para los cuales el entrelazamiento algebraico no
resuelve por sí solo los problemas de dominio.

\begin{theorem}[Límite inductivo acotado]
\label{pf84:E03-06}
Sean \(J_m:\mathcal H_m\to\mathcal H_{m+1}\) isometrías y
\(A_m\in\mathcal B(\mathcal H_m)\) operadores autoadjuntos tales que
\[
 A_{m+1}J_m=J_mA_m,
 \qquad
 \sup_m\lVert A_m\rVert<\infty.
\]
Existe un único operador acotado autoadjunto \(A\) sobre
\(\mathcal H_\infty=\varinjlim(\mathcal H_m,J_m)\) cuya restricción a cada
nivel es \(A_m\).
\end{theorem}

\begin{proof}
En la unión algebraica densa se define
\[
 A[J_{m,\infty}\xi]:=J_{m,\infty}A_m\xi.
\]
El entrelazamiento hace que la definición no dependa del representante. La
cota uniforme permite extender \(A\) por continuidad a
\(\mathcal H_\infty\); la simetría en el subespacio denso y la acotación
implican su autoadjunción.
\end{proof}

\begin{remark}[Control para operadores no acotados]
\label{pf84:E03-07}
Si los \(A_m\) no son acotados, la identidad
\(A_{m+1}J_m=J_mA_m\) sólo define una acción en una unión de dominios.
Para obtener un operador cerrado o autoadjunto deben declararse los dominios
transportados, probarse cerrabilidad y controlarse los índices de
deficiencia, o bien demostrarse esencial autoadjunción sobre un núcleo común
compatible. Omitir estas obligaciones constituye un fallo de tipo, no una
abreviación de la prueba.
\end{remark}

Sea ahora \(\mathcal S_\infty\) un conjunto no vacío de secciones
supervivientes compatibles y sea \(\mathcal P_m\) el conjunto finito de sus
prefijos de longitud \(m\). Escribimos
\[
 \mathcal H_m:=\ell^2(\mathcal P_m),
 \qquad
 E_m(\gamma,\gamma')=
 \begin{cases}
 1,&\gamma'\text{ prolonga }\gamma,\\
 0,&\text{en otro caso},
 \end{cases}
\]
y
\[
 d_m(\gamma):=
 \#\{\gamma'\in\mathcal P_{m+1}:\gamma'\mapsto\gamma\},
 \qquad
 D_m^{\deg}:=E_mE_m^*=\operatorname{diag}(d_m(\gamma)).
\]
Nos restringimos al subespacio vivo \(d_m(\gamma)>0\), y cada prefijo de
nivel \(m+1\) posee un único padre.

\begin{theorem}[Isometría normalizada de prolongación]
\label{pf84:C01}
El operador
\[
 J_m:=E_m^*(D_m^{\deg})^{-1/2}:
 \mathcal H_m\longrightarrow\mathcal H_{m+1}
\]
es una isometría, y
\[
 J_me_\gamma=
 \frac1{\sqrt{d_m(\gamma)}}
 \sum_{\gamma'\mapsto\gamma}e_{\gamma'}.
\]
\end{theorem}

\begin{proof}
La unicidad del padre da
\[
 J_m^*J_m
 =(D_m^{\deg})^{-1/2}E_mE_m^*(D_m^{\deg})^{-1/2}=I.
\]
La fórmula sobre \(e_\gamma\) se obtiene desarrollando la columna
correspondiente de \(E_m^*\).
\end{proof}

Puesto
\[
 \mathcal K_m:=\mathcal H_{m+1}\ominus J_m\mathcal H_m,
 \qquad
 \mathcal H_{\mathrm{surv}}
 \cong\mathcal H_0\oplus\bigoplus_{m\ge0}\mathcal K_m,
\]
se define el operador de refinamiento
\[
 \mathcal R|_{\mathcal H_0}=0,
 \qquad
 \mathcal R|_{\mathcal K_m}=3^{m+1}I_{\mathcal K_m},
\]
con dominio maximal
\[
 \Dom(\mathcal R)=
 \mathcal H_0\oplus
 \left\{(\xi_m)_{m\ge0}:
 \sum_{m\ge0}3^{2m+2}\lVert\xi_m\rVert^2<\infty\right\}.
\]

\begin{theorem}[Operador de refinamiento]
\label{pf84:C02}
Si cada \(\mathcal K_m\) es finito dimensional, \(\mathcal R\) es positivo,
autoadjunto y de resolvente compacto. Su espectro mide profundidad de
refinamiento; no constituye un espectro de primos ni de ceros.
\end{theorem}

\begin{proof}
La descomposición ortogonal reduce \(\mathcal R\) a bloques escalares reales
sobre su dominio maximal, lo que prueba positividad y autoadjunción. Los
autovalores \(3^{m+1}\) escapan a infinito y sus multiplicidades son finitas;
por tanto, el resolvente es compacto. Un bloque infinito dimensional a
autovalor acotado falsaría esta última conclusión.
\end{proof}

\section{Operadores entero, primo y orbital}

Sea \(\mathscr W_{\mathrm{bal}}\) el conjunto de palabras ternarias
balanceadas reducido mediante el teorema
\cref{pf84:C03}. Sobre \(\ell^2(\mathscr W_{\mathrm{bal}})\), el operador
\[
 Be_w=b(w)e_w
\]
se toma en su dominio maximal. Su restricción al subespacio de los enteros
positivos se denota por \(Q\).

\begin{theorem}[Operadores entero y positivo]
\label{pf84:C04}
El operador \(B\) es autoadjunto, posee espectro simple \(\ZZ\) y
resolvente compacto. El operador \(Q\) posee espectro simple
\(\NN_{\ge1}\) y resolvente compacto.
\end{theorem}

\begin{proof}
La biyección de \cref{pf84:C03} enumera cada entero exactamente una vez. Una
diagonal real sobre su dominio maximal es autoadjunta; como sus entradas
escapan a infinito fuera de conjuntos finitos, su resolvente es compacto. La
restricción positiva conserva estas propiedades.
\end{proof}

La prueba clásica omitida de inversión theta, continuación y ecuación
funcional se fusiona una sola vez con la realización triádica ya construida
en \cref{pf84:C05}. Allí se fijan literalmente
\(\Theta_B=1+2\Theta_+\) y \(Z=\mathcal Z_3\); por tanto, este bloque no
define un segundo origen de \(\zeta\).

Sea \(\mathbb P\) el conjunto de primos ordinarios,
\(\mathfrak h_{\mathbb P}:=\ell^2(\mathbb P)\) y
\[
 H_{\mathbb P}e_p=(\log p)e_p.
\]
En la Fock bosónica \(\Gamma_s(\mathfrak h_{\mathbb P})\), las ocupaciones
\((k_p)_p\) tienen soporte finito.

\begin{theorem}[Representación de Fock de la factorización]
\label{pf84:C08}
La aplicación sobre bases
\[
 U_F\lvert(k_p)_p\rangle=e_{\prod_pp^{k_p}}
\]
se extiende a un unitario
\[
 U_F:\Gamma_s(\mathfrak h_{\mathbb P})
 \longrightarrow\ell^2(\NN_{\ge1})
\]
y satisface
\[
 U_F\,d\Gamma(H_{\mathbb P})\,U_F^*=\log Q.
\]
\end{theorem}

\begin{proof}
La factorización única da una biyección entre ocupaciones finitas y enteros
positivos, luego una identificación unitaria de las bases ortonormales.
Además,
\[
 \sum_pk_p\log p=\log\!\left(\prod_pp^{k_p}\right),
\]
lo que prueba el entrelazamiento. La estructura de partida explícita es la
elección de los modos primos y de la Fock simétrica.
\end{proof}

Sobre
\[
 \mathcal H_{\mathrm{orb}}
 :=\ell^2\{(p,k,\varepsilon):
 p\in\mathbb P,\ k\ge1,\ \varepsilon=\pm1\},
\]
definimos
\[
 D_{\mathrm{orb}}e_{p,k,\varepsilon}
 =\varepsilon k\log p\,e_{p,k,\varepsilon},
 \qquad
 W_{\log}e_{p,k,\varepsilon}
 =(\log p)e_{p,k,\varepsilon}.
\]

\begin{proposition}[Traza orbital]
\label{pf84:C09}
El operador \(D_{\mathrm{orb}}\) es autoadjunto y de resolvente compacto. Si
\(f\in C_c(\RR)\) es acotada, entonces
\[
 \Tr\!\left(
 W_{\log}e^{-|D_{\mathrm{orb}}|/2}f(D_{\mathrm{orb}})
 \right)
 =
 \sum_{p,k}(\log p)p^{-k/2}
 \bigl[f(k\log p)+f(-k\log p)\bigr].
\]
\end{proposition}

\begin{proof}
La diagonal es real y sus longitudes escapan de todo compacto con
multiplicidad finita. El soporte compacto de \(f\) hace que el producto
regulado tenga rango finito; su traza es la suma de las entradas diagonales.
El índice \(p\) recorre aquí los primos ordinarios. El resultado realiza el
lado primo clásico, pero no identifica \(D_{\mathrm{orb}}\) con un operador
dual de los ceros.
\end{proof}

\section{Forma completada de Guinand--Weil}

Fijamos
\[
 \mathcal D_{\mathrm W}:=C_c^\infty(\RR),
 \qquad
 \widehat\psi(t)=\int_{\RR}\psi(x)e^{-itx}\,\dd x,
 \qquad
 \psi(x)=\frac1{2\pi}\int_{\RR}\widehat\psi(t)e^{itx}\,\dd t,
\]
y, para \(\psi,\phi\in\mathcal D_{\mathrm W}\),
\[
 h_{\psi,\phi}(z)
 :=\widehat\psi(z)\,
 \overline{\widehat\phi(\overline z)}.
\]
En la recta real \(h_{\psi,\psi}=|\widehat\psi|^2\).

\begin{definition}[Clase de prueba]
\label{pf84:E01}
La clase de prueba es \(\mathcal D_{\mathrm W}\), estable por derivación,
reflexión y traslación. Se utiliza la inversa
\[
 \check h(x)=\frac1{2\pi}\int_{\RR}h(t)e^{itx}\,\dd t.
\]
La suavidad y el soporte compacto proporcionan decaimiento rápido en bandas
y soporte compacto para \(\check h\), condiciones que garantizan la
convergencia de los bloques siguientes.
\end{definition}

Sea \(\psi_\Gamma=\Gamma'/\Gamma\). Definimos
\[
 \Irr_\#:=
 \left\{\operatorname{ev}_9(w):
 w\in\mathscr W_9^+,\ 
 P_{\Irr,\#}e_w=e_w\right\}.
\]
Por \cref{pf84:C07}, esta imagen coincide con los primos ordinarios; aquí se
usa únicamente para tipar el índice del bloque orbital, sin repetir aquella
prueba.
\begin{align}
 \mathcal P(h)&:=h(i/2)+h(-i/2),\\
 \mathcal G(h)&:=-\log\pi\,\check h(0)
 +\frac1{2\pi}\int_{\RR}h(t)
 \Re\psi_\Gamma\!\left(\frac14+\frac{it}{2}\right)\dd t,\\
 \mathcal O_\#(\check h)&:=
 \sum_{p\in\Irr_\#}\sum_{k\ge1}
 (\log p)p^{-k/2}
 \bigl[\check h(k\log p)+\check h(-k\log p)\bigr].
\end{align}

\begin{definition}[Funcional completado]
\label{pf84:E02}
Para \(h=h_{\psi,\phi}\), se fija
\[
 \mathcal W(h):=\mathcal P(h)+\mathcal G(h)
 -\mathcal O_\#(\check h).
\]
El bloque orbital procede del selector irreducible del
\cref{chap:producto-nativo}; los bloques polar y gamma pertenecen a la
completación clásica de zeta. Esta diferencia de procedencia se conserva en
la notación y en la prueba.
\end{definition}

\begin{theorem}[Fórmula explícita de Guinand--Weil]
\label{pf84:E03}
Para \(\psi,\phi\in C_c^\infty(\RR)\) y \(h=h_{\psi,\phi}\),
\[
 \lim_{T\to\infty}
 \sum_{\substack{\rho:\,\zeta(\rho)=0\\|\Im\rho|\le T}}
 h\!\left(\frac{\rho-1/2}{i}\right)
 =\mathcal W(h),
\]
donde los ceros no triviales se cuentan con multiplicidad y el límite se
toma con la simetrización clásica.
\end{theorem}

\begin{proof}
Paley--Wiener coloca \(h_{\psi,\phi}\) en la clase admisible de la fórmula
explícita. Los polos de
\(\Lambda_\zeta(s)=\pi^{-s/2}\Gamma(s/2)\zeta(s)\) producen
\(\mathcal P\), el factor arquimediano produce \(\mathcal G\) y el producto
de Euler produce la suma sobre potencias primas. La identificación de esta
última suma con \(\mathcal O_\#\) se realiza únicamente mediante
\cref{pf84:E04}; así no se duplica la prueba de primalidad nativa.
\end{proof}

\begin{proposition}[Reindexación del bloque orbital]
\label{pf84:E04}
Si los irreducibles del producto nativo coinciden exactamente con los primos
ordinarios y se conservan los pesos
\((\log p)p^{-k/2}\) y las dos orientaciones \(\pm k\log p\), entonces
\[
 \mathcal O_\#(\check h)=\mathcal O(\check h).
\]
\end{proposition}

\begin{proof}
La biyección entre irreducibles nativos y primos reindexa la suma sin
modificar índices, pesos ni argumentos. Un irreducible no primo, un primo
ausente o un cambio de peso falsaría la identidad.
\end{proof}

\begin{remark}[Separación meromorfa--entera]
\label{pf84:E05}
La función \(\Lambda_\zeta\) posee polos simples en \(0\) y \(1\), que
originan el bloque polar. En cambio,
\[
 \xi(s)=\frac12s(s-1)\Lambda_\zeta(s)
\]
es entera: los factores \(s(s-1)\) cancelan precisamente esos polos. No se
atribuyen, por tanto, polos en \(0\) o \(1\) a \(\xi\).
\end{remark}

\begin{definition}[Forma hermitiana y preoperador]
\label{pf84:E06}
Se definen
\[
 \mathcal B_{\mathrm W}(\psi,\phi)
 :=\mathcal W(h_{\psi,\phi}),
 \qquad
 \mathcal V_{\mathrm W}:=
 \mathcal D_{\mathrm W}/\operatorname{Rad}\mathcal B_{\mathrm W},
\]
y, sobre el cociente algebraico, la regla candidata
\[
 D_{\mathrm W,0}[\psi]:=[-i\psi'].
\]
No se presupone que \(\mathcal B_{\mathrm W}\) sea positiva ni que
\(\mathcal V_{\mathrm W}\) sea ya un espacio de Hilbert. La regla sólo queda
bien definida después de probar la invariancia del radical.
\end{definition}

\begin{theorem}[Simetría algebraica e invariancias]
\label{pf84:E07}
El radical es invariante bajo \(-i\,\dd/\dd x\);
\(D_{\mathrm W,0}\) está bien definido y es
\(\mathcal B_{\mathrm W}\)-simétrico. Las traslaciones
\(T_a\psi(x)=\psi(x-a)\) preservan la forma y la reflexión
\(J\psi(x)=\psi(-x)\) satisface
\[
 JD_{\mathrm W,0}J=-D_{\mathrm W,0}.
\]
\end{theorem}

\begin{proof}
La convención de Fourier da
\[
 h_{-i\psi',\phi}(z)=h_{\psi,-i\phi'}(z).
\]
La identidad vale antes de aplicar cualquiera de los tres bloques de
\(\mathcal W\), de modo que la derivada preserva el radical y el
preoperador es simétrico. Las fases opuestas de las traslaciones se cancelan
en \(h\); la reflexión cambia el signo de la derivada y deja invariantes los
bloques polar, gamma y orbital. Esta simetría algebraica no equivale todavía
a autoadjunción en un Hilbert positivo.
\end{proof}

\section{Geometría local de la recta crítica}

\begin{lemma}[Coordenada de Möbius de la recta crítica]
\label{pf84:RH-01}
Para \(s\in\CC\setminus\{0,1\}\), sea
\[
 \tau(s):=i\frac{s-1}{s}.
\]
Con \(F(s)=1-s\) y \(C(s)=\overline s\),
\[
 \tau(Fs)=-\frac1{\tau(s)},
 \qquad
 \tau(Cs)=-\overline{\tau(s)},
\]
y
\[
 \tau(Fs)=\tau(Cs)
 \iff |\tau(s)|=1
 \iff \Re s=\frac12.
\]
\end{lemma}

\begin{proof}
Las dos identidades se obtienen por sustitución directa. Su igualdad
equivale a \(\tau(s)^{-1}=\overline{\tau(s)}\), es decir,
\(|\tau(s)|=1\); desarrollar esta última condición da
\(\Re s=1/2\). Para los ceros de zeta se requieren además las simetrías
funcional y conjugada. La equivalencia parametriza la recta; no sitúa por sí
sola los ceros en ella.
\end{proof}

\begin{lemma}[Signo local de Hermite--Biehler]
\label{pf84:RH-04}
Sea \(A\) una función real entera con un cero real simple \(a\). Para
\(z=a+iy\), \(y>0\) suficientemente pequeño,
\[
|A(z)-iA'(z)|^2-|A(z)+iA'(z)|^2
 =-4A'(a)^2y+O(y^2)<0.
\]
La convención local asociada a \(A\) es
\[
 E_A:=A+iA'.
\]
Sólo al especializar \(A=\Xi\) y suponer además que \(a\) es un cero real
simple de \(\Xi\) se obtiene
\[
 E_\Xi=\Xi+i\Xi'.
\]
\end{lemma}

\begin{proof}
Se desarrolla \(A\) y \(A'\) en Taylor alrededor de \(a\), usando
\(A(a)=0\) y \(A'(a)\ne0\), y se retienen los términos de primer orden en
\(y\). El resultado es estrictamente local: no demuestra que \(E_A\) posea la
propiedad global de Hermite--Biehler.
\end{proof}

\begin{proposition}[Obstrucción por leyes de conteo]
\label{pf84:E14}
Sea
\[
 N_A(T):=\Tr\mathbf1_{[-T,T]}(A).
\]
Los operadores \(Q\), \(B\), \(D_{\mathrm{orb}}\) y \(\mathcal R\) no
poseen, por sí solos, la ley de conteo exigida a un eventual operador dual
de los ceros.
\end{proposition}

\begin{proof}
Para \(T\ge1\),
\[
 N_Q(T)=\lfloor T\rfloor,
 \qquad
 N_B(T)=2\lfloor T\rfloor+1.
\]
El teorema de los números primos da
\[
 N_{D_{\mathrm{orb}}}(T)
 =2\sum_{p\le e^T}\left\lfloor\frac{T}{\log p}\right\rfloor
 \ge2\pi(e^T)\sim\frac{2e^T}{T}.
\]
Finalmente, \(N_{\mathcal R}(T)\) es constante en cada intervalo
multiplicativo \([3^m,3^{m+1})\), pues \(\mathcal R\) sólo tiene niveles
\(3^{m+1}\) con multiplicidades finitas.

Para
\(\Xi(z)=\xi(1/2+iz)\), la ecuación funcional implica
\(\Xi(-z)=\Xi(z)\); los ceros \(\pm\gamma\) tienen la misma multiplicidad y
\[
 N_{\mathrm{sym}}(T)=2N_+(T)+O(1).
\]
En consecuencia, la convención bilateral requiere
\[
 N_{\mathrm{sym}}(T)\sim\frac{T}{\pi}\log T,
\]
mientras la unilateral equivalente, bajo esta simetría y emparejamiento de
multiplicidades, requiere
\[
 N_+(T)\sim\frac{T}{2\pi}\log T.
\]
Ninguna de las cuatro funciones anteriores tiene esa asintótica.
\end{proof}

\section{Laplacianos sobre los ciclos de orden \texorpdfstring{$3^m$}{potencias de tres}}

Para $m\geq 1$, sean

\[
 N_m=3^m,
 \qquad
 C_m=\ZZ/N_m\ZZ,
\]

y considérese el laplaciano combinatorio

\[
 (\Delta_m f)(j)=2f(j)-f(j+1)-f(j-1),
 \qquad f\in\ell^2(C_m).
\]

Los caracteres

\[
 e_{m,n}(j)=N_m^{-1/2}
 \exp\!\left(\frac{2\pi i n j}{N_m}\right),
 \qquad n\in C_m,
\]

forman una base ortonormal y satisfacen

\[
 \Delta_m e_{m,n}
 =4\sin^2\!\left(\frac{\pi n}{N_m}\right)e_{m,n}.
\]

Como $N_m$ es impar, cada clase de $C_m$ posee un único representante
$\widetilde n$ con

\[
 -\frac{N_m-1}{2}\leq\widetilde n\leq\frac{N_m-1}{2}.
\]

Definimos el operador autoadjunto $D_m$ por

\begin{equation}
 D_m e_{m,n}
 =\frac{N_m}{\sqrt\pi}
 \sin\!\left(\frac{\pi\widetilde n}{N_m}\right)e_{m,n}.
 \label{eq:operador-espectral-Dm}
\end{equation}

Entonces

\begin{equation}
 D_m^2=\frac{N_m^2}{4\pi}\,\Delta_m.
 \label{eq:Dm-laplaciano}
\end{equation}

Sea $P_0$ la proyección ortogonal sobre el carácter constante. Para toda
función par acotada $F$ sobre el espectro de $D_m$, introducimos la traza
positiva

\begin{equation}
 \operatorname{Tr}^{+}_m\!\bigl(F(D_m)\bigr)
 :=\frac12\operatorname{Tr}
 \!\left(F(D_m)P_0^{\perp}\right).
 \label{eq:traza-positiva}
\end{equation}

El factor $1/2$ selecciona una contribución por cada par espectral
$\{n,-n\}$; la eliminación de $P_0$ excluye el autovalor nulo.

\begin{theorem}[Límite theta de las trazas de calor]
\label{thm:limite-theta-triadico}
Para todo $x>0$,
\[
 \lim_{m\to\infty}
 \operatorname{Tr}^{+}_m\!\left(e^{-xD_m^2}\right)
 =\Theta_+(x)
 :=\sum_{n\geq1}e^{-\pi n^2x}.
\]
\end{theorem}

\begin{proof}
La traza bilateral, una vez retirado el modo constante, es
\[
 \sum_{\substack{n\in C_m\\n\neq0}}
 \exp\!\left[
 -\frac{N_m^2x}{\pi}
 \sin^2\!\left(\frac{\pi\widetilde n}{N_m}\right)
 \right].
\]
Para cada entero fijo $n$, el sumando converge a $e^{-\pi n^2x}$. Además,
para $|\widetilde n|\leq N_m/2$, la desigualdad
$\sin y\geq 2y/\pi$, válida en $0\leq y\leq\pi/2$, proporciona una
mayorante gaussiana sumable e independiente de $m$. La convergencia
dominada aplicada a la suma discreta implica
\[
 \lim_{m\to\infty}
 \operatorname{Tr}\!\left(e^{-xD_m^2}P_0^\perp\right)
 =\sum_{n\in\ZZ\setminus\{0\}}e^{-\pi n^2x}.
\]
La división entre dos establecida en \eqref{eq:traza-positiva} conserva una
contribución de cada par $\{n,-n\}$ y da $\Theta_+(x)$.
\end{proof}

\section{Transformada theta--Mellin}

Para $\Re s>1$, definimos

\begin{equation}
 \mathcal Z_{3}(s)
 :=
 \frac{\displaystyle
 \int_0^\infty\Theta_+(x)x^{s/2-1}\,\dd x}
 {\pi^{-s/2}\Gamma(s/2)}.
 \label{eq:funcional-theta-mellin}
\end{equation}

La convergencia absoluta permite integrar término a término. Como

\[
 \int_0^\infty e^{-\pi n^2x}x^{s/2-1}\,\dd x
 =\Gamma(s/2)(\pi n^2)^{-s/2},
\]

se obtiene

\begin{equation}
 \mathcal Z_3(s)=\sum_{n\geq1}n^{-s}=\zeta(s),
 \qquad \Re s>1.
 \label{eq:Z3-zeta}
\end{equation}

La igualdad se prolonga meromórficamente por la continuación de la función
zeta de Riemann \cite{Apostol1976}. En lo sucesivo, todas las evaluaciones
proceden de \eqref{eq:funcional-theta-mellin}, de sus derivadas o de la
descomposición residual de su serie de Dirichlet.

\begin{theorem}[Inversión theta y ecuación funcional]
\label{pf84:C05}
Con la notación operatorial precedente,
\[
 \Theta_B(x):=\Tr(e^{-\pi xB^2})
 =1+2\Theta_+(x),
 \qquad
 Z(s):=\Tr(Q^{-s})=\mathcal Z_3(s)
 \quad(\Re s>1).
\]
La sumación de Poisson y la transformada de Mellin producen
\[
 \Theta_B(x)=x^{-1/2}\Theta_B(1/x)
\]
y la continuación de
\[
 \xi(s)=\frac12s(s-1)\pi^{-s/2}\Gamma(s/2)Z(s),
 \qquad
 \xi(s)=\xi(1-s).
\]
\end{theorem}

\begin{proof}
La igualdad \(\Theta_B=1+2\Theta_+\) separa el modo cero y empareja
\(\pm n\); la identidad \(Z=\mathcal Z_3\) es
\eqref{eq:Z3-zeta}. La sumación de Poisson aplicada a la gaussiana da la
inversión de \(\Theta_B\). En el semiplano de convergencia, la integración
término a término proporciona
\[
 \pi^{-s/2}\Gamma(s/2)Z(s)
 =\int_0^\infty x^{s/2-1}\Theta_+(x)\,\dd x.
\]
Al dividir la integral en \((0,1)\) y \((1,\infty)\), aplicar la inversión
theta al primer tramo y separar los dos polos clásicos, se obtiene la
continuación y la ecuación funcional. Esta única prueba clásica incorporada
no localiza los ceros de \(\zeta\).
\end{proof}

\section{Descomposición residual módulo tres}

Para $r\in\{0,1,2\}$, sea

\begin{equation}
 Z_r(s)=
 \sum_{\substack{n\geq1\\n\equiv r\pmod3}}n^{-s}.
 \label{eq:canales-residuales-Zr}
\end{equation}

En términos de la zeta de Hurwitz,

\[
 Z_0(s)=3^{-s}\zeta(s),
 \qquad
 Z_1(s)=3^{-s}\zeta\!\left(s,\frac13\right),
 \qquad
 Z_2(s)=3^{-s}\zeta\!\left(s,\frac23\right).
\]

Los tres términos tienen residuo $1/3$ en $s=1$. Escribimos sus
desarrollos de Laurent en la forma

\begin{equation}
 Z_r(1+z)-\frac1{3z}=\sum_{n\geq0}c_{r,n}\,z^n,
 \qquad
 \zeta(1+z)-\frac1z=\sum_{n\geq0}a_n\,z^n.
 \label{eq:coeficientes-Laurent-residuales}
\end{equation}

\begin{theorem}[Descomposición de los coeficientes de Laurent]
\label{thm:descomposicion-laurent-triadica}
Para todo $n\geq0$,
\[
 a_n=c_{0,n}+c_{1,n}+c_{2,n},
 \qquad
 \gamma_n=(-1)^n n!\,a_n.
\]
Si $\chi_{-3}$ denota el carácter real no principal módulo tres, entonces
\[
 \frac{L^{(n)}(1,\chi_{-3})}{n!}=c_{1,n}-c_{2,n}.
\]
Además, con \(\lambda=\log3\),
\[
 c_{0,n}=\frac13\left[
 \frac{(-\lambda)^{n+1}}{(n+1)!}
 +\sum_{k=0}^{n}a_k\frac{(-\lambda)^{n-k}}{(n-k)!}
 \right].
\]
Estas tres identidades determinan de manera única
\((c_{0,n},c_{1,n},c_{2,n})\) a partir de
\((a_n,L^{(n)}(1,\chi_{-3}))\) y de los coeficientes agregados de orden
menor.
\end{theorem}

\begin{proof}
La partición de los enteros positivos por clases residuales proporciona
$\zeta=Z_0+Z_1+Z_2$. Al sustraer los polos en $s=1$ y comparar los
coeficientes de \eqref{eq:coeficientes-Laurent-residuales} se obtiene la
primera identidad. La segunda coincide con la convención usual para las
constantes de Stieltjes. Además,
$Z_1-Z_2=L(s,\chi_{-3})$; los polos de $Z_1$ y $Z_2$ se cancelan y la
comparación de coeficientes da la tercera relación. Por último,
\[
 Z_0(1+z)=\frac13e^{-\lambda z}\zeta(1+z).
\]
El producto de las dos series de Laurent proporciona la fórmula explícita
para \(c_{0,n}\). La suma y la diferencia recuperan entonces
\(c_{1,n}\) y \(c_{2,n}\), lo que prueba la unicidad afirmada.
\end{proof}

La suma de los tres términos recupera el sector agregado; la diferencia de
las clases $1$ y $2$ conserva la orientación del carácter; la comparación
con la clase divisible por tres registra el cambio de escala. Estas tres
combinaciones son linealmente independientes y determinan los tres
coeficientes $c_{r,n}$.

\section{Euler--Mascheroni y constantes de Stieltjes}

Sea

\[
 C_r:=\operatorname*{FP}_{s=1}Z_r(s).
\]

Las identidades clásicas para la función digamma dan
\cite{Apostol1976,Lagarias2013}

\begin{equation}
 C_0=\frac{\gamma-\log3}{3},
 \label{eq:C0-parte-finita}
\end{equation}

\begin{equation}
 C_1=\frac\gamma3+\frac{\log3}{6}+\frac{\pi}{6\sqrt3},
 \qquad
 C_2=\frac\gamma3+\frac{\log3}{6}-\frac{\pi}{6\sqrt3}.
 \label{eq:C1-C2-partes-finitas}
\end{equation}

\begin{corollary}
\label{cor:gamma-log3-L1}
Las partes finitas satisfacen
\[
 \boxed{\gamma=C_0+C_1+C_2},
\]
\[
 \boxed{L(1,\chi_{-3})=C_1-C_2=\frac{\pi}{3\sqrt3}},
\]
\[
 \boxed{\log3=C_1+C_2-2C_0}.
\]
\end{corollary}

\begin{proof}
En orden cero, el \cref{thm:descomposicion-laurent-triadica} da
$a_0=c_{0,0}+c_{1,0}+c_{2,0}$. Por definición,
$a_0=\gamma$ y $c_{r,0}=C_r$, lo que prueba la primera igualdad. La
diferencia de las dos clases orientadas es
$C_1-C_2=L(1,\chi_{-3})$. Sustituyendo
\eqref{eq:C0-parte-finita}--\eqref{eq:C1-C2-partes-finitas} se obtienen las
otras dos identidades.
\end{proof}

Una única extracción de coeficientes mediante la fórmula integral de
Cauchy aplicada a \eqref{eq:coeficientes-Laurent-residuales} determina
simultáneamente $\gamma_0,\ldots,\gamma_6$. Las discretizaciones con 96 y
192 nodos coinciden hasta orden $10^{-87}$; la identidad de los
coeficientes, no obstante, procede del
\cref{thm:descomposicion-laurent-triadica} y no de esa comparación numérica.

\section{Evaluación de la constante de Apéry}

\begin{proposition}
\label{prop:evaluacion-apery}
El funcional \eqref{eq:funcional-theta-mellin} satisface
\[
 \boxed{
 \mathcal Z_3(3)=\zeta(3)
 =1.202056903159594285399738161511449\ldots}.
\]
\end{proposition}

\begin{proof}
La igualdad exacta es la especialización $s=3$ de \eqref{eq:Z3-zeta}.
Para certificar cualquier prefijo decimal sin modificar el funcional,
considérese $S_N=\sum_{n=1}^{N}n^{-3}$. El criterio integral da
\[
 S_N+\frac1{2(N+1)^2}
 <\zeta(3)<
 S_N+\frac1{2N^2}.
\]
Los intervalos racionales anidados resultantes determinan las cifras
impresas.
\end{proof}

La irracionalidad de $\zeta(3)$ es el teorema clásico de Apéry
\cite{Apery1979}; no se deduce de la mera evaluación del funcional.

\section{Constante de Glaisher--Kinkelin}

La derivada espectral en $s=-1$ define

\begin{equation}
 \boxed{
 A_{\mathrm{GK}}
 =\exp\!\left(\frac1{12}-\mathcal Z_3'(-1)\right)}.
 \label{eq:glaisher-kinkelin-espectral}
\end{equation}

Esta es la fórmula del determinante regularizado asociado al producto
hiperfactorial \cite{Vardi1988}. Una derivación central de cinco puntos,
evaluada a dos escalas, proporciona

\[
 A_{\mathrm{GK}}
 =1.2824271291006226368753425688697917\ldots.
\]

El término $1/12$ pertenece a la regularización del hiperfactorial y no se
introduce como parámetro de ajuste.

\section{Constante de Euler--Kronecker del campo de Eisenstein}

El carácter $\chi_{-3}$ determina el campo cuadrático imaginario

\[
 F_{-3}=\QQ(\sqrt{-3}),
 \qquad
 \zeta_{F_{-3}}(s)=\zeta(s)L(s,\chi_{-3}).
\]

Su constante de Euler--Kronecker es

\begin{equation}
 \boxed{
 \gamma_{F_{-3}}
 =\gamma+\frac{L'(1,\chi_{-3})}{L(1,\chi_{-3})}}.
 \label{eq:euler-kronecker-eisenstein}
\end{equation}

Los coeficientes residuales de
\eqref{eq:coeficientes-Laurent-residuales} dan

\[
 L(1,\chi_{-3})
 =0.6045997880780726168646927525473852\ldots,
\]
\[
 L'(1,\chi_{-3})
 =0.2226629869686015094866602627647443\ldots,
\]

y, por \eqref{eq:euler-kronecker-eisenstein},

\[
 \gamma_{F_{-3}}
 =0.9454972808716807032397499941581890\ldots.
\]

La elección de $F_{-3}$ procede del carácter residual módulo tres; no de
una búsqueda sobre los valores decimales de constantes de campos
cuadráticos.

\section{Producto local asociado a los primos gemelos}

Para un patrón finito $H=\{h_1,\ldots,h_k\}\subset\ZZ$, definimos

\[
 \nu_p(H):=\left|\{-h_1,\ldots,-h_k\}\bmod p\right|,
 \qquad
 \sigma_p(H)=
 \frac{1-\nu_p(H)/p}{(1-1/p)^k}.
\]

Para $H=\{0,2\}$, los factores locales conducen a

\begin{equation}
 C_2^{\mathrm{twin}}
 =\prod_{p>2}\frac{p(p-2)}{(p-1)^2}.
 \label{eq:constante-primos-gemelos}
\end{equation}

Sea \(P_X\) el producto restringido a los primos \(2<p\leq X\). Para
\(p>X\), todos los factores pertenecen a \((0,1)\); al ampliar el producto
de la cola desde los primos a todos los enteros se obtiene la cota
telescópica
\[
 \prod_{n>X}\left(1-\frac1{(n-1)^2}\right)
 =\frac{X-1}{X}
 \leq
 \prod_{p>X}\left(1-\frac1{(p-1)^2}\right)<1.
\]
El producto \(P_{2\times10^6}\), evaluado con redondeo decimal dirigido, y
esta desigualdad racional para la cola dan

\begin{equation}
 0.6601615071224244826403933551
 <C_2^{\mathrm{twin}}<
 0.6601618372035081248537566591.
 \label{eq:intervalo-primos-gemelos}
\end{equation}

El recuento directo hasta $10^6$ contiene 8169 pares de primos gemelos. La
convergencia del producto local y este censo finito no implican la
existencia de infinitos pares; esta última es una proposición aritmética
distinta.

\section{Constante de Meissel--Mertens y Hamiltoniano primo}

La constante de Meissel--Mertens se define por

\begin{equation}
 B_1=\lim_{x\to\infty}
 \left(\sum_{p\leq x}\frac1p-\log\log x\right)
 \label{eq:meissel-mertens-def}
\end{equation}

y satisface

\begin{equation}
 B_1=\gamma+
 \sum_p\left(\log\!\left(1-\frac1p\right)+\frac1p\right).
 \label{eq:meissel-mertens-euler}
\end{equation}

Con la notación de \cref{pf84:C08}, ponemos
\[
 \mathcal F_P:=\Gamma_s(\mathfrak h_{\mathbb P}),
 \qquad
 H_P:=d\Gamma(H_{\mathbb P}).
\]
No se redefine aquí la Fock bosónica de modos primos ni su Hamiltoniano. El
operador ya construido tiene función de partición

\begin{equation}
 \operatorname{Tr}_{\mathcal F_P}(e^{-sH_P})
 =\prod_p\frac1{1-p^{-s}}
 =\zeta(s),
 \qquad \Re s>1.
 \label{eq:hamiltoniano-primo-zeta}
\end{equation}

Sea \(\mathcal H_P^{(1)}\) el subespacio de una partícula, generado por
\(\{|p\rangle:p\text{ primo}\}\). Entonces
\[
 \operatorname{Tr}_{\mathcal H_P^{(1)}}\!\left(
 e^{-H_P}\mathbf 1_{(-\infty,\log x]}(H_P)\right)
 =\sum_{p\leq x}\frac1p.
\]
Por tanto, \(B_1\) es la parte finita, cuando \(x\to\infty\), de esta traza
espectral truncada después de sustraer \(\log\log x\). La igualdad
\eqref{eq:hamiltoniano-primo-zeta} presupone la descomposición de Fock sobre
los primos; por ello constituye una representación operatorial del producto
de Euler, no una deducción de la primalidad a partir del laplaciano
triádico.

\section{Dependencia respecto del operador}

Las evaluaciones anteriores comparten el operador
\eqref{eq:operador-espectral-Dm}, la traza positiva
\eqref{eq:traza-positiva} y la transformada
\eqref{eq:funcional-theta-mellin}. Otras constantes requieren dinámicas
distintas. La constante de Catalan está asociada al carácter módulo cuatro
$\chi_{-4}$, mientras que la constante de Khinchin pertenece al operador de
Gauss y a su medida invariante. Un prefijo de cualquiera de ellas puede
aparecer en un catálogo finito sin que el funcional triádico la seleccione.
Esta dependencia operatorial distingue la evaluación espectral de la
capacidad puramente combinatoria de codificación.
